\section[拓扑量子力学、拓扑 Maxwell 理论与类域论]{拓扑量子力学、拓扑 Maxwell 理论与类域论（2 月 11 日）}

\srcpage{24}

上次我们稍微谈到，在拓扑量子力学中，我们在导出的意义下杀死时间演化算子：先加一个微分 $Q$，使 Hilbert 空间成为链复形（chain complex），然后说明 $H$ 恰当（exact），或说同伦于 $0$。你随后可以拿上同调当作你的 Hilbert 空间；但如今人们已经意识到，更方便的做法是在导出意义下保留整个链复形而令 $H$ 恰当。

分次（由 $\Z$ 或 $\Z/2$ 给出）是超对称代数作用结构的一部分，但从某种意义上说，谁在先？也许你是在通过 R-对称性研究该代数的表示（这里就是 $\mathfrak{u}_1$，分次由此而来），分次随后才出现；或者你先有分次，然后造出代数作用。

实现这件事有两种不同方式。

（1）在「A 型拓扑量子力学」中，令 $\mathcal{H} := \Omega^\bullet_X$、$Q := d$、$Q^* := d^*$，以及 $H := \Delta$，即 Hodge Laplacian。这就是我们上次讨论的例子。

（2）在「B 型拓扑量子力学」中，我们不用 Riemann 流形，而令 $X$ 是复流形。那么附近有一个类似的复形（你可以叫它复复形）：$\mathcal{H} := \Omega^{0,\bullet}_X$、$Q := \bar\partial$、$Q^* := \bar\partial^*$，以及 $H := \Delta_{\bar\partial}$。关于 $Q$ 的上同调是 $H^{0,*}(X)$，它拟同构于 $R\Gamma(\mathcal{O}_X)$，即 $X$ 上函数层的导出整体截面。

这两者都是一维 $\mathcal{N} = 1$ 超对称代数的表示。这里 $\mathcal{N}$ 度量「有多少超对称」；有精确的办法定义它，但眼下最好把它理解为：这告诉你一维中可能的超对称量是最小的。

还有 $\mathcal{N} = 2$ 超对称代数，以及 $\mathcal{N} = 2$ 版本的拓扑量子力学。这里我们把上面两种类型结合起来：令 $X$ 是紧 Kähler 流形，链复形双分次为 $\Omega^{\bullet,\bullet}_X$，而我们需要更多算子：$\partial$、$\partial^*$、$\bar\partial$ 与 $\bar\partial^*$。有一个深刻的数学事实（硬 Lefschetz 定理）：$\mathrm{SL}_2(\C)$ 作用在上同调上；对角部分给出分次，此外还有提升与降低算子。R-对称群是 $\mathrm{SU}_2\subset\mathrm{SL}_2(\C)$。

从数学观点看这是个困难而深刻的定理，但在物理中它以更初等的方式出现：你问 Kähler 流形上量子力学的超对称代数是什么，而一个表示论的事实（它未必更简单）告诉你我们得到 $\mathcal{N} = 2$，从而有 $\mathrm{SL}_2(\C)$-作用，因为这本来就出现在 $d = 1$ 的 $\mathcal{N} = 2$ 超对称代数中。

还可以再推一点：若 $X$ 是紧超 Kähler（hyperKähler）流形，超对称代数扩大为 $\mathcal{N} = 4$，而上同调不再有 $\mathrm{SL}_2(\C)$-作用，而是 $\mathrm{Spin}_5$-作用，作为 $\mathcal{N} = 4$ 超 Lie 代数的一部分。你甚至可以一路走到 $\mathcal{N} = 8$。这些代数由线性无关的奇算子（超荷，supercharge）的个数区分：$\mathcal{N} = 1$ 中两个，$\mathcal{N} = 2$ 中四个，$\mathcal{N} = 4$ 中八个，$\mathcal{N} = 8$ 中十六个。奇次算子空间的维数是思考 $\mathcal{N}$ 的更好的数学方式，不过你也看得出这并不是 $\mathcal{N}$ 的字面含义。

往这些理论里再加一些场，我们会发现更多超对称，即扩大作用在理论上的超 Lie 群。我们的要求（desiderata）是 Hamilton 量 $H = [Q,-]$，即它是恰当的。而对拓扑理论我们想做得更多：量子力学中一般有一个算子 $T$ 称为应力张量（stress tensor），度量对度规的依赖，我们希望它也 $Q$-恰当，从而让理论更拓扑：在上同调层面什么都不依赖度规。Costello--Gwilliam 的书里对此有更详细的解释。

\textbf{拓扑 Maxwell 理论}\quad 现在把这件事煮到 $4$ 维：取 Maxwell 理论，添加场，并找出超对称，以便在这个 $Q$-恰当的意义下得到拓扑场论。这是一个 $\mathcal{N} = 4$ 理论，有十六个超荷和一大堆数据，包括一个带联络的复线丛、
\srcpage{25}
一个称为 Higgs 场（Higgs field）的 $1$-形式，一个复标量 $u$，以及四个费米子（它们给出奇算子）。

同样有两个选项，以 A 与 B 标号。对 A-扭（A-twist），我们 $3$-流形 $M$ 上的 Hilbert 空间是联络空间的上同调，其中有些微妙之处。首先，联络是一个栈（stack），因为有些联络带有稳定化群的作用；其次规范群在作用，我们取等变上同调。也就是说，$H = H^\bullet_{U_1}(\mathcal{C}(M))$。所以这就像在 $U_1$-联络空间上做 A 型拓扑量子力学（de Rham 理论）。

另一条路是 B-扭：我们令 $\nabla$ 是主 $\C^\times$-丛上的联络，而不是 $U_1$-丛上的。然后我们在 $\C^\times$-联络空间上做 B 型拓扑量子力学（Dolbeault 那套故事）。

回到 A 侧片刻。我们会稍微滥用记号，说 $\mathrm{Pic}(M)$ 就是拓扑空间 $\mathcal{C}(M)$。于是 $H^0(\mathrm{Pic}(M))$ 是 $\mathcal{C}(M)$ 上的局部常值函数。$\mathcal{C}(M)$ 的连通分支是主 $U_1$-丛的拓扑类型，即 $H^2(M;\Z)$。所以前面几讲我们想过的东西又回来了，只是语境更花哨。

对更高阶的上同调，我们把 $\mathrm{Pic}(M)$ 分解为格 $\times$ 环面 $\times$ 向量空间 $\times\,BU_1$，这是我们需要考虑的。我们看到 $H^2(M;\Z)$、环面 $H^1(M;U_1)$，以及来自 $BU_1$ 上同调的 $\C[u]$。在拓扑扭之下，事情简单得多。

在 B 侧，我们关心 $\mathrm{Loc}_{\C^\times}(M)$，它看起来像复环面 $T^\vee_\C$，即映射 $\pi_1(M)\to\C^\times$ 的环面。这些映射通过 $H_1(M)$ 分解，所以这是对偶于格 $H_1(M)$ 的环面。于是我们的 Hilbert 空间看起来像 $\C[H^2(M;\Z)]$ 或 $\C[T^\vee_\C]$，而有 Poincaré--Pontrjagin 对偶把它们互换。

一个收获是：我们先前看到的、关于流形上同调的 Fourier 级数，作为拓扑 Maxwell 理论故事的一部分出现了；此外还有新的东西在发生。A-扭的超 Maxwell 理论研究联络空间的拓扑，B-扭版本研究平坦联络空间的代数几何。电磁对偶把两者互换。这两个理论不是同一理论的两种描述，而是两个真正不同的理论之间的对偶。

A 侧 $\mathrm{Pic}$ 的拓扑对偶于 $\mathrm{Loc}$ 的代数几何，这个视角我们会一再看到。

幕后说一句：这两个拓扑场论是某个超对称 QFT——称为 $\mathcal{N} = 4$ 超 Maxwell 理论——的拓扑扭。它是我们这个世界里光的理论的超对称版本。这里的 $\mathcal{N} = 4$ 告诉我们面对的是复几何（很像一维中 $\mathcal{N} = 2$ 是 $X$ 为 Kähler 的信号）。关于 $\mathcal{N} = 4$ 超 Yang--Mills 的文献很多；这是规范群为 $U_1$ 的特例，所以若要找参考文献，去查超 Yang--Mills。

对偶再一次互换 Wilson 算子与 't Hooft 算子。在 A 侧我们有 't Hooft 算子，它在格上平移；在 B 侧我们有 Wilson 算子，它乘以沿曲线 $\gamma\in H_1(M;\Z)$ 的单值（monodromy）。't Hooft 算子造出磁单极子。物理理论同时有电与磁，但选取这两种扭中的任一种，我们就破坏了这个对称性：在 A 侧只有电的测量，在 B 侧只有磁的测量。

现在来引入缺陷（defect）。它们会稍微改变场空间。

\begin{itemize}
  \item 首先是「类时」't Hooft 算子。取点 $x\in M$，考虑在 $x$ 处坐着一个单极子时的电磁学。也就是说，在 $M\times[0,1]$ 中，直线 $\{x\}\times[0,1]$ 是单极子的世界线：它不动。随后我们考虑这个单极子存在时的电磁学，办法是考虑 $M\setminus x$ 上的联络。这个空间分裂为不交并
  \[
    \mathcal{C}(M\setminus x) = \coprod_{c_1(P|_{S_\varepsilon(x)}) = n}\mathcal{C}(M,nx),
  \]
  即按 $x$ 的环绕球面上的 Chern 类取指标。我们把 $n$ 看成单极子的荷。于是我们得到这个联络空间上 $L^2$ 函数构成的 Hilbert 空间。

  \item 在 B 侧，我们要 $\mathrm{Loc}_{\C^\times}(M,x)$ 上的「类时」Wilson 算子；它定义为平坦 $\C^\times$-联络连同 $x$ 处纤维的一个平凡化。这是 $\mathrm{Loc}_{\C^\times}(M)$ 上的主 $\C^\times$-丛。这个「额外的 $\C^\times$」对偶于 A 侧荷的「额外的 $\Z$」，而 Fourier 级数把这两个额外群上的函数空间等同起来。物理上这被解释为造出一个带电粒子，而不必为这个粒子添加新的场。
\end{itemize}

所以电磁对偶告诉我们：造一个磁单极子，对偶于造一个重到足以破坏此对称性的带电粒子（我们并没有添加对应于这个粒子的新场）。

\srcpage{26}

还有更多缺陷要考虑。下一个称为分歧（ramification），或说加一个面缺陷（surface defect）或螺线管（solenoid）；它对应物理中如下情形：考虑一股电流流过紧紧绕在柱体上、形成一个管子的导线。我们的场在这里会做出古怪的事；具体地说，我们从 $M$ 中挖掉一条曲线 $\beta$，考虑 $\mathrm{Loc}_{\C^\times}(M\setminus\beta)$。此时有 $\pi_1$，它在现实世界中意味着 Aharonov--Bohm 效应：若你绕着螺线管走一圈，电场会受影响。（TODO：也许我 Arun 误解了这一点。）

于是我们可以看这个局部系统空间上的全纯函数（这是 B 侧，即电侧）。在 A 侧，即磁侧，我们看 $M$ 上、沿 $\beta$ 有平凡化的联络空间的上同调。这是一个映射 $\beta\to U_1 = S^1$，它有绕数，从而添加一个 $\Z$，而在对偶侧我们曾有一个 $\C^\times$（绕额外环路的和乐）。同样，$\Z$ 与 $\C^\times$ 上的函数通过 Fourier 变换等同。

这里的数学全是对 $3$-流形上 Poincaré 对偶的解释，而这些 Fourier 理论的陈述各自容许不同的物理含义，这很妙。

我们这里的物理故事是关于几何 Langlands 的；对应的算术故事是什么？这是数论中称为类域论（class field theory）的学科。同样，我们只会得到交换 $\Z$ 与 $\C^\times$（更一般地，格与环面）的 Fourier 理论，但我们会在有趣的地方、以有趣的方式找到这些群。同样地，某个叫 $\mathrm{Pic}$ 的东西的拓扑与某个叫 $\mathrm{Loc}$ 的东西的代数几何互换。起初，「拓扑」只意味着 $\pi_0$，而「代数几何」意味着全纯函数，但可以提取更多结构。

现在说设定。设 $F$ 是 $\Q$ 的有限扩张。数论学家把重要的对象附着到 $F$ 上，即（理想）类群 $\mathrm{Cl}(F) := \Pic(\Spec\mathcal{O}_F)$。也就是说，它告诉我们 $\Spec\mathcal{O}_F$ 上的线丛，即秩 $1$ 投射 $\mathcal{O}_F$-模，同构类下的情形。不过在这个语境中它还有另一个更初等的解释：$\mathcal{O}_F$ 的分式理想（即落在 $F$ 中者）模主理想。类群总是有限的，这相当漂亮，而在其他场合没有明显的类比。数域最重要的不变量之一是类数 $n_F := |\mathrm{Cl}(F)|$。

（无分歧）类域论把类群与 Galois 理论联系起来。它把 $\C[\mathrm{Cl}(F)]$ 与 $\C[\mathrm{Loc}_{\mathbb{G}_m}(\mathcal{O}_F)]$ 等同。你若去翻类域论的教科书，会发现记号很不一样，但这样做是为了与今天早些时候我们那个更几何、更物理的故事对接。那么 $\mathrm{Loc}$ 到底是什么？它是 $\pi_1$-表示，而我们的 $\pi_1$ 的类比是 Galois 群。对流形而言，$\pi_1$ 是万有覆叠在底空间上的自同构；在我们的算术拓扑词典下，它对应于域扩张 $F^{\mathrm{un}}/F$ 的自同构，其中 $F^{\mathrm{un}}$ 是 $F$ 的最大无分歧阿贝尔扩张。

为什么是它而不是 $F$？我们不把 $F$ 想成一个点 $\Spec F$，而想成概形 $\Spec(\mathcal{O}_F)$ 上的有理函数。代数闭包对应于允许带 punctures 的曲线上的函数域，但眼下不必操心这个，先想闭曲线。

\begin{definition}
  域扩张 $F\hookrightarrow E$ 在素理想 $\mathfrak{p}\in\mathcal{O}_F$ 处无分歧（unramified），是指 $\mathcal{O}_E\otimes_{\mathcal{O}_F}\mathcal{O}_F/\mathfrak{p}$ 没有幂零元。
\end{definition}
\yizhu{底本此处标为 8.2，而按本工程体例，条目编号由 LaTeX 按讲次重新起算，故这里显示为 8.1；号差源于底本自身的编号空缺。}{B（体例说明，原文 p.26）}

这些幂零算子对应于分歧覆盖中的分歧点：这里有些有趣的几何在发生。无论如何，无分歧阿贝尔扩张构成偏序集，它有最大元，那就是 $F^{\mathrm{un}}$。而这就是 $\Spec\mathcal{O}_F$ 的 étale 基本群。

这里有个技术条件：我们还要求事物在无穷位处也无分歧。取 $F\otimes_\Q\R$ 会怎样？它同构于 $\R^r\times\C^s$，其中 $r$ 与 $s$ 数的是 $F$ 到某些东西里的嵌入。我们说 $E/F$ 在无穷处无分歧，是指这里没有扩张 $\R\hookrightarrow\C$。例如，全虚域（即 $r = 0$ 者）在无穷处无分歧。

无论如何，无分歧类域论的主定理是 $\mathrm{Cl}(F)\cong\Gal(F^{\mathrm{un}}/F)$。这是 Hilbert 类域 $F^{\mathrm{un}}$ 的自同构群。我们想用 $\mathrm{Pic}$ 与 $\mathrm{Loc}$ 来解释它。类群是 Picard 群；另一侧把 $\Hom(\Gal(F^{\mathrm{un}}/F),\C^\times)$ 解释为 $H^1(\Spec(\mathcal{O}_F))$（TODO：怎么解释？）。Galois 群的特征标对偶也是 $\Hom(\mathrm{Cl}(F),\C^\times) = \mathrm{Cl}(F)^\vee$。

这个语境中有 Poincaré 对偶的类比，称为 Artin--Verdier 对偶，它适用于整数环的谱。这是算术拓扑词典的一部分：数域像 $3$-流形的又一种表现。Artin--Verdier 对偶的证明用到类域论，所以这与我们在拓扑情形中的论证大体上是倒过来的。

\srcpage{27}

另外，$\Spec(\mathcal{O}_F)$ 更像无定向的 $3$-流形：有某个东西给出对偶，但它不是常值层。取而代之，我们用分圆特征标（cyclotomic character）得到某个东西。这与 $\Z/n$ 和 $n$ 次单位根 $\mu_n$ 的差别有关，也与 $\Q/\Z$ 和全体单位根的差别、$\C^\times$ 和 $\mathbb{G}_m$ 的差别有关。对偶化对象不是 $\Z$，而是 Tate 扭 $\Z(1)$；回到物理里可以把它解释为：对应着把 Fourier 变换中出现的 $2\pi i$ 因子算进去。这些细节使 Artin--Verdier 对偶的精确陈述更令人困惑，但关键思想是 Poincaré 对偶的类比。

下次我们会讨论这个故事的函数域版本，以及别的一些东西。
